\chapter{결과의 재현}\label{app:repro}

이 보고서의 수치는 저장소 루트에서 다음 명령으로 다시 만들 수 있다. 측정 결과의 원본은
\code{docs/results/}에 있다.

\begin{table}[ht]
\centering
\small
\begin{tabularx}{\linewidth}{>{\raggedright\arraybackslash}p{50mm}>{\raggedright\arraybackslash}X}
\toprule
결과 & 명령 \\
\midrule
거짓 정지의 통계 (그림~\ref{fig:stopstats}, \ref{fig:nodiag}) &
  \code{cargo run -{}-release -p christopher-core -{}-example stop\_statistics 400} \\
전체 탐색 성능 (표~\ref{tab:perf}) &
  \code{cargo run -{}-release -p christopher-core -{}-example full\_search enigma-i} (또는 \code{m4}) \\
예제 전문과 실제 전문의 해독 (\ref{ch:software}장, \ref{sec:historical}절) &
  노트북 \code{notebooks/01\_crib\_to\_plaintext.ipynb}를 처음부터 실행 \\
기준 RTL의 시험 (표~\ref{tab:rtl-tests}) & \code{rtl/tb/run.sh} (Verilator와 cocotb가 필요하다) \\
\bottomrule
\end{tabularx}
\end{table}

그림~\ref{fig:stopstats}와 \ref{fig:nodiag}의 자료는 \code{report/data/}의 CSV 파일이며, 각 파일의 머리에
출처를 적어 두었다. 보고서는 \code{report/} 폴더에서 \code{latexmk}로 빌드한다.

\chapter{배선표}\label{app:wirings}

에니그마 로터 I--VIII, 그리스 로터 Beta와 Gamma, 반사판 네 종의 배선이다. 이 표는 손으로 옮겨 적지 않았다.
저장소의 배선 원본 \code{data/wirings/enigma.txt} 하나를 Rust 코어, RTL 패키지(\code{rtl/gen/gen\_wirings.py}),
드럼 PCB 생성기(\code{hardware/gear/drum-pcb/gen\_drum.py})가 모두 읽고, 이 부록의 표도 같은 파일에서
\code{report/figures/make\_appendix\_tables.py}가 만든다. 생성 스크립트는 표를 쓰기 전에 각 배선이 순열인지,
반사판이 고정점 없는 대합인지 다시 검사한다. 배선이 맞다는 근거는 실제 전문의 해독(부록~\ref{app:vectors})과
\code{AAAAA}$\to$\code{BDZGO} 같은 알려진 결과다.

배선은 코어 오프셋 0에서 \emph{입구 쪽 글자 $\to$ 반사판 쪽 글자}다. 코어가 $s$만큼 돌아 있으면 글자 $c$는
$W[c+s]-s$로 나간다(식~\ref{eq:through}). 같은 배선이 기어식 드럼 PCB의 두 배선 벌(가는 길, 오는 길)과 릴레이식의
로터 배선 카드에 그대로 쓰인다. 오는 길에는 역배선이 쓰인다.

% GENERATED by report/figures/make_appendix_tables.py from data/wirings/enigma.txt — do not edit.
\begin{table}[ht]\centering\small
\caption{로터와 그리스 로터의 배선(코어 오프셋 0). 입력 글자 A--Z 아래에 출력 글자를 적었다. 역배선은 반사판 쪽에서 들어오는 오는 길이다. 넘김 글자는 그 글자가 창에 보일 때 다음 키에서 왼쪽 로터를 함께 돌린다.}
\label{tab:app-rotors}
\begin{tabular}{@{}l l l c@{}}\toprule
로터 & 배선 (입력 A--Z) & 역배선 & 넘김\\\midrule
입력 & \texttt{ABCDEFGHIJKLMNOPQRSTUVWXYZ} & \texttt{ABCDEFGHIJKLMNOPQRSTUVWXYZ} & \\\midrule
I & \texttt{EKMFLGDQVZNTOWYHXUSPAIBRCJ} & \texttt{UWYGADFPVZBECKMTHXSLRINQOJ} & Q\\
II & \texttt{AJDKSIRUXBLHWTMCQGZNPYFVOE} & \texttt{AJPCZWRLFBDKOTYUQGENHXMIVS} & E\\
III & \texttt{BDFHJLCPRTXVZNYEIWGAKMUSQO} & \texttt{TAGBPCSDQEUFVNZHYIXJWLRKOM} & V\\
IV & \texttt{ESOVPZJAYQUIRHXLNFTGKDCMWB} & \texttt{HZWVARTNLGUPXQCEJMBSKDYOIF} & J\\
V & \texttt{VZBRGITYUPSDNHLXAWMJQOFECK} & \texttt{QCYLXWENFTZOSMVJUDKGIARPHB} & Z\\
VI & \texttt{JPGVOUMFYQBENHZRDKASXLICTW} & \texttt{SKXQLHCNWARVGMEBJPTYFDZUIO} & ZM\\
VII & \texttt{NZJHGRCXMYSWBOUFAIVLPEKQDT} & \texttt{QMGYVPEDRCWTIANUXFKZOSLHJB} & ZM\\
VIII & \texttt{FKQHTLXOCBJSPDZRAMEWNIUYGV} & \texttt{QJINSAYDVKBFRUHMCPLEWZTGXO} & ZM\\
Beta & \texttt{LEYJVCNIXWPBQMDRTAKZGFUHOS} & \texttt{RLFOBVUXHDSANGYKMPZQWEJICT} & ---\\
Gamma & \texttt{FSOKANUERHMBTIYCWLQPZXVGJD} & \texttt{ELPZHAXJNYDRKFCTSIBMGWQVOU} & ---\\
\bottomrule\end{tabular}
\end{table}

\begin{table}[ht]\centering\small
\caption{반사판의 배선. 반사판은 글자를 둘씩 맞바꾸는 대합(involution)이고 고정점이 없다. 그래서 에니그마는 어떤 글자도 자기 자신으로 암호화하지 않는다(크립 정렬의 근거, \ref{sec:crib}절). 얇은 반사판(B-THIN, C-THIN)은 M4에서 그리스 로터와 함께 쓴다.}
\label{tab:app-reflectors}
\begin{tabular}{@{}l l l@{}}\toprule
반사판 & 배선 (입력 A--Z) & 맞바뀌는 쌍\\\midrule
B & \texttt{YRUHQSLDPXNGOKMIEBFZCWVJAT} & \texttt{\footnotesize AY BR CU DH EQ FS GL IP JX KN MO TZ VW}\\
C & \texttt{FVPJIAOYEDRZXWGCTKUQSBNMHL} & \texttt{\footnotesize AF BV CP DJ EI GO HY KR LZ MX NW QT SU}\\
B-THIN & \texttt{ENKQAUYWJICOPBLMDXZVFTHRGS} & \texttt{\footnotesize AE BN CK DQ FU GY HW IJ LO MP RX SZ TV}\\
C-THIN & \texttt{RDOBJNTKVEHMLFCWZAXGYIPSUQ} & \texttt{\footnotesize AR BD CO EJ FN GT HK IV LM PW QZ SX UY}\\
\bottomrule\end{tabular}
\end{table}


\chapter{검증 벡터와 조작 절차}\label{app:vectors}

\section{검증 벡터}\label{sec:app-vectors}

검증 벡터는 두 종류다. \emph{실제 전문 벡터}는 에니그마 구현이 역사적 기계와 같은지를 확인한다(표~\ref{tab:app-historical}).
\emph{회로 검증 벡터}는 RTL과 FPGA 엔진이 소프트웨어 봄베와 같은 답을 내는지를 확인한다(표~\ref{tab:app-rtl}).
두 표 모두 저장소의 파일에서 \code{report/figures/make\_appendix\_tables.py}가 생성했다.

회로 검증 벡터는 모두 Rust 코어가 낸 값이고, 회로는 이것을 그대로 재현해야 한다. 한 위치 벡터는 그 위치에서 676선이
도달한 고정점 전체(\code{live}: 케이블마다 26비트)와 정지 여부를 담는다. 그래서 전파 회로의 사소한 차이도 정지 여부가
같은 경우까지 포함해 잡힌다. 전체 구동 벡터는 로터 순서 하나를 끝까지 돈 정지 목록이고, 회로는 정지의 \emph{순서}까지
같아야 한다(드럼 순회 순서가 같으므로). 조건이 다른 경우(대각선 보드 없음, 가운데 로터 넘김 가설, M4)를 일부러 섞었다.
벡터의 형식과 다시 만드는 방법은 \code{data/vectors/rtl/README.md}에 있다.

% GENERATED by report/figures/make_appendix_tables.py from data/vectors — do not edit.
\begin{table}[ht]\centering\footnotesize
\caption{실제 전문 벡터(\code{data/vectors/historical.json}). 링과 시작 위치는 왼쪽부터(M4는 그리스 로터가 맨 앞). 확인 열은 저장소의 시험(\code{crates/christopher-core/tests/historical.rs})이 무엇을 대조하는지다.}
\label{tab:app-historical}
\begin{tabular}{@{}p{3.4cm} l p{2.2cm} l r p{3.4cm}@{}}\toprule
벡터 & 기종 & 로터 순서 & 링 / 시작 & 길이 & 확인\\\midrule
\texttt{m4-\allowbreak{}triton-\allowbreak{}shark-\allowbreak{}1945-\allowbreak{}03-\allowbreak{}06} & m4 & B-thin Gamma IV-V-III & \texttt{AAPU} / \texttt{ZKSA} & 344 & 공개 평문과 전체 일치\\
\texttt{barbarossa-\allowbreak{}heeresgruppe-\allowbreak{}nord-\allowbreak{}1941-\allowbreak{}07-\allowbreak{}08-\allowbreak{}nr23} & enigma-i & B IV-III-II & \texttt{PKF} / \texttt{PIK} & 120 & 공개 평문 없음; 해독문이 독일어 통계로 암호문보다 뚜렷이 높음\\
\bottomrule\end{tabular}
\end{table}

\paragraph{\code{m4-triton-shark-1945-03-06}.} 플러그보드 \texttt{AX BZ CY DG ET IP JN OQ RV SW}. 출처: \url{https://cryptocellar.org/bgac/m4-messages.html}.
\begin{quote}\footnotesize\ttfamily
암호문(앞 50자): NEWME SUCII FPRLE ETPNZ LPYJW DKLLD UJURJ LBUKF WNKRS PENZU\\
평문(앞 50자): EULAU FENNA CHPER IODEJ BRUNO JXXEI NSKKF UNKFE UERGE UPPEE
\end{quote}

\paragraph{\code{barbarossa-heeresgruppe-nord-1941-07-08-nr23}.} 플러그보드 \texttt{CY EL FH GS IJ KQ MW PV RZ TU}. 출처: \url{https://cryptocellar.org/bgac/heeresgruppe-nord-july-1941.pdf}.
\begin{quote}\footnotesize\ttfamily
암호문(앞 50자): CWSEB DDIRB ZUUBG KJANB VGIVD VDZZI GAKBZ CJMMV EMVXL THNLD\\
\end{quote}

\begin{table}[ht]\centering\footnotesize
\caption{회로 검증 벡터(\code{data/vectors/rtl/}). 위: \code{bombe\_core} 한 위치씩의 벡터(676선 고정점 전체 \code{live}와 정지 여부). 아래: \code{bombe\_sweep}·FPGA 엔진이 순서까지 맞춰야 하는 전체 구동의 정지 목록. 모두 Rust 코어(\code{Bombe.trace}, \code{Bombe.run})가 낸 값이다.}
\label{tab:app-rtl}
\begin{tabular}{@{}p{4.6cm} l r r r p{2.4cm}@{}}\toprule
파일 & 로터 순서 & 연결 & 벡터 & 정지 & 조건\\\midrule
\texttt{enigma-\allowbreak{}i\_\allowbreak{}B-\allowbreak{}III-\allowbreak{}V-\allowbreak{}II.jsonl} & B III-V-II & 14 & 65 & 1 & ---\\
\texttt{enigma-\allowbreak{}i\_\allowbreak{}nodiag\_\allowbreak{}B-\allowbreak{}II-\allowbreak{}IV-\allowbreak{}V.jsonl} & B II-IV-V & 13 & 63 & 23 & 대각선 보드 없음\\
\texttt{enigma-\allowbreak{}i\_\allowbreak{}turnover5\_\allowbreak{}B-\allowbreak{}I-\allowbreak{}II-\allowbreak{}III.jsonl} & B I-II-III & 11 & 80 & 40 & 넘김 가설 5\\
\texttt{enigma-\allowbreak{}i\_\allowbreak{}weak\_\allowbreak{}C-\allowbreak{}III-\allowbreak{}V-\allowbreak{}II.jsonl} & C III-V-II & 8 & 80 & 45 & ---\\
\texttt{m4\_\allowbreak{}B-\allowbreak{}THIN\_\allowbreak{}BETA\_\allowbreak{}II-\allowbreak{}IV-\allowbreak{}I.jsonl} & B-THIN BETA II-IV-I & 11 & 100 & 40 & ---\\
\texttt{m4\_\allowbreak{}C-\allowbreak{}THIN\_\allowbreak{}GAMMA\_\allowbreak{}VI-\allowbreak{}VIII-\allowbreak{}VII\_\allowbreak{}turnover9.jsonl} & C-THIN GAMMA VI-VIII-VII & 11 & 80 & 40 & 넘김 가설 9\\
\midrule
파일 & 로터 순서 & 연결 & 위치 & 정지 & 조건\\\midrule
\texttt{sweep\_\allowbreak{}enigma-\allowbreak{}i\_\allowbreak{}B-\allowbreak{}III-\allowbreak{}V-\allowbreak{}II\_\allowbreak{}11edges.json} & B III-V-II & 11 & 17,576 & 82 & ---\\
\texttt{sweep\_\allowbreak{}enigma-\allowbreak{}i\_\allowbreak{}B-\allowbreak{}III-\allowbreak{}V-\allowbreak{}II\_\allowbreak{}14edges.json} & B III-V-II & 14 & 17,576 & 1 & ---\\
\texttt{sweep\_\allowbreak{}enigma-\allowbreak{}i\_\allowbreak{}C-\allowbreak{}V-\allowbreak{}IV-\allowbreak{}I\_\allowbreak{}13edges\_\allowbreak{}nodiag\_\allowbreak{}turnover5.json} & C V-IV-I & 13 & 17,576 & 16 & 대각선 보드 없음, 넘김 가설 5\\
\texttt{sweep\_\allowbreak{}m4\_\allowbreak{}B-\allowbreak{}THIN\_\allowbreak{}BETA\_\allowbreak{}II-\allowbreak{}IV-\allowbreak{}I\_\allowbreak{}11edges.json} & B-THIN BETA II-IV-I & 11 & 456,976 & 2,187 & ---\\
\bottomrule\end{tabular}
\end{table}


\section{조작 절차}\label{sec:app-procedure}

\begin{todobox}{조작 절차}
Python 인터페이스, FPGA 호스트 프로그램, 실물의 조작 순서.
\end{todobox}
